\section{Integral generators and the dimension squeeze}
\label{sec:induction}

We now combine the upper bound on the full equation space with the
odd-weight values constructed by holonomy.  The integral choice of
generators is part of the argument: a rational depth-one value need
not initially have the desired reduction modulo two.

\subsection{Depth and integral independence}
For a Lie polynomial in \(x,y\), its ordinary \(y\)-depth is the
smallest number of occurrences of \(y\) in any nonzero associative
word.  Write \(\mathcal D^p L\) for the subspace of depth at least
\(p\).  A homogeneous Lie polynomial of weight greater than one
has depth at least one, since the free Lie algebra on \(x\) alone
has no such homogeneous component.
This depth is distinct from the \(B\)-count used after reduction
modulo two: \(C=[x,y]\) contributes one \(y\) but has \(B\)-count zero.

\begin{lemma}\label{lem:depth-addition}
If \(\psi\in\mathcal D^p L\) and
\(\phi\in\mathcal D^q L\), then
\[
 \{\psi,\phi\}\in\mathcal D^{p+q}L.
\]
Consequently every Lie word of length at least two in homogeneous
elements of weights greater than one has zero depth-one component.
\end{lemma}
\begin{proof}
An application of \(D_\psi\) replaces one occurrence of \(y\) by
\([y,\psi]\), while it kills \(x\).  Each nonzero resulting word
therefore has at least \(p\) more occurrences of \(y\) than the
original word.  Thus
\(D_\psi(\phi)\in\mathcal D^{p+q}L\).
The same reasoning applies to \(D_\phi(\psi)\), and the ordinary
bracket also adds depth.  Equation~\eqref{eq:ihara} proves the
assertion.  Its consequence follows by induction on the length of
the Lie word.
\end{proof}

\begin{lemma}\label{lem:hall-saturation}
Let \(K\) be a finite free \(\Z_{(2)}\)-module, and suppose that
\(w_1,\ldots,w_s\in K\) have linearly independent reductions in
\(K/2K\).  Then they are linearly independent over \(\Q\), and
\[
 \left(\sum_{i=1}^s\Q w_i\right)\cap K
   =\sum_{i=1}^s\Z_{(2)}w_i.
\]
\end{lemma}
\begin{proof}
In a nonzero rational relation, multiply by a rational scalar so
that every coefficient belongs to \(\Z_{(2)}\) and at least one is
a unit.  Reduction modulo two gives a contradiction.

Now let \(v=\sum_i c_iw_i\in K\), with \(c_i\in\Q\).
If some \(c_i\notin\Z_{(2)}\), let \(e\ge1\) be the smallest integer
for which every \(2^ec_i\) belongs to \(\Z_{(2)}\).
At least one of these coefficients is a unit, whereas
\[
 \sum_i(2^ec_i)w_i=2^ev\in2K.
\]
Reduction gives the same contradiction.  Thus all the \(c_i\)
belong to \(\Z_{(2)}\), proving the equality.
\end{proof}

Let \(V=\bigoplus_{n\ge2}V_n\) be the rational subalgebra of
categorical values from Proposition~\ref{prop:categorical-values},
so that \(V_n\subset W_n\).
Let \(K_n\) be the ordinary free-Lie \(\Z_{(2)}\)-lattice of
weight \(n\), and put
\[
 M_n=V_n\cap K_n,\qquad L_n=W_n\cap K_n.
\]
By Proposition~\ref{prop:saturated-reduction}, their reductions
\(\overline M_n\subset\overline L_n\subset K_n/2K_n\) have
dimensions \(\dim_\Q V_n\) and \(\dim_\Q W_n\), respectively.
The Ihara bracket has integral coefficients on Lie words.
Since \(V\) is a Lie subalgebra, \(M=\bigoplus_n M_n\) is closed
under that bracket.

Recall the free Lie algebra
\[
 \mathfrak f_{\mathrm{odd}}
   =\Lie_{\F}\langle z_k:k\ge1\rangle,\qquad
 \wt(z_k)=2k+1,
\]
and let \(d_{n,r}\) and \(d_n=\sum_r d_{n,r}\) be its dimensions
in weight \(n\) and length \(r\), and in weight \(n\), respectively.
These dimensions agree with their characteristic-zero counterparts
by Lemma~\ref{lem:lie-lattice}.
Corollary~\ref{cor:dimension-bound} gives
\begin{equation}\label{eq:dimension-squeeze}
 \dim_\Q V_n
 \le \dim_\Q W_n
 =\dim_\F\overline L_n
 \le\sum_r d_{n,r}=d_n.
\end{equation}
All sums here are finite.  In particular, we have not assumed
that every characteristic-two solution lifts to a rational one.

\subsection{Choosing the next generator}
We construct \(\sigma_{2k+1}\in M_{2k+1}\) inductively so that
its reduction has leading \(B\)-count one, with projection
\begin{equation}\label{eq:generator-leading-term}
 \pi\bigl(\operatorname{in}_1\overline\sigma_{2k+1}\bigr)
       =\ad_C^kB.
\end{equation}
Here \(\operatorname{in}_1\) denotes the count-one component and
\(\pi\) sets \(A=0\).  We identify the right side with \(z_k\).

Suppose the generators of all odd weights below \(n\) have been
chosen.  Choose a Hall basis for the free Lie algebra on their
formal symbols, and consider its words of weight \(n\).
All intermediate values lie in \(M\), since \(M\) is closed under
the Ihara bracket.
If such a word has length \(r\), repeated application of
Proposition~\ref{prop:leading-ihara} places its reduced value in
filtration \(F^r\) and identifies its projected leading part with
the same Hall word in the \(z_k\)'s.
Those leading parts are linearly independent for each \(r\).
They are independent across different \(r\) as well: in any
putative relation take the smallest length occurring, pass to
that associated graded piece, and project.  Its coefficients
must all vanish; repeat for the remaining lengths.

Thus the Hall-word values have independent reductions in
\(K_n/2K_n\).  Lemma~\ref{lem:hall-saturation} proves both their
rational independence and the fact that their rational span
intersects \(K_n\) in precisely their \(\Z_{(2)}\)-span.
Each of these Hall words has length at least two, since the
available generators have weights below \(n\).  Every integral
vector in their rational span therefore still reduces into
\(F^2\), and cannot provide the required leading-count-one class.

Every free Lie word of weight \(n\) and length at least two uses
only generators of weights below \(n\).  The only possible
missing Hall word is the single generator of weight \(n\), which
exists exactly when \(n\) is odd and \(n\ge3\).  Therefore the old
Hall-word values supply
\[
 d_n-\varepsilon_n
 \quad\text{independent elements of }V_n,\qquad
 \varepsilon_n=
 \begin{cases}
  1,&n\ge3\text{ odd},\\
  0,&n\text{ even}.
 \end{cases}
\]

If \(n\) is even, this already attains the upper bound
\eqref{eq:dimension-squeeze}.  If \(n\ge3\) is odd,
Proposition~\ref{prop:odd-depth-one} gives
\(\tau_n\in V_n\) with a nonzero coefficient of \(x^{n-1}y\).
Every old Hall-word value of weight \(n\) is decomposable and
has zero depth-one component by Lemma~\ref{lem:depth-addition}.
Hence \(\tau_n\) is independent of them.  In both cases,
\begin{equation}\label{eq:dimension-equality}
 \dim_\Q V_n=\dim_\Q W_n=d_n.
\end{equation}

For odd \(n=2k+1\), we must still obtain the integral choice
\eqref{eq:generator-leading-term}.
The filtration on \(\overline M_n\) is induced by its embedding in
the ambient ordinary Lie algebra on \(A,C,B\).  Its associated
graded embeds in that of \(\overline L_n\), and
Corollary~\ref{cor:dimension-bound} gives
\[
 \dim_\F\gr^r\overline M_n\le d_{n,r}.
\]
Equation~\eqref{eq:dimension-equality} and saturation give equality
of the sums of these dimensions.  Each individual deficit is
nonnegative, so equality holds for every \(r\).
For \(r=1\), the target is the one-dimensional line spanned by
\(z_k\).  The injective projection is therefore an isomorphism
onto that line.
Choose its nonzero preimage in \(\gr^1\overline M_n\), represent
it by a vector of \(F^1\overline M_n\), and lift that vector to
\(M_n\).  This lift is \(\sigma_n\).  Its projected coefficient
is one because the coefficient field is \(\F\).
There is no requirement that the initially constructed
\(\tau_n\) have this reduction.

The induction begins at \(n=3\), with no previously chosen
generators.  The same argument supplies \(\sigma_3\).
Weights one and two vanish by Lemma~\ref{lem:low-weights}.
Proceeding through successive weights proves
\eqref{eq:dimension-equality} in all weights and makes every
required choice.

\subsection{Freeness, spanning, and completion}
\begin{proof}[Proof of Theorem~\ref{thm:main}]
\label{thm:main-proof}
Send \(e_{2k+1}\) to the elements \(\sigma_{2k+1}\) just
constructed.  Since the target \(V\) is a Lie algebra, the
universal property gives a graded Lie homomorphism
\[
 \Lie_\Q\langle e_3,e_5,\ldots\rangle\longrightarrow V.
\]
In each weight, its Hall-basis images have independent
reductions by the preceding leading-term argument.  They are
therefore rationally independent, so the map is injective.
Their number in weight \(n\) is \(d_n\), and
\eqref{eq:dimension-equality} shows that they span both \(V_n\)
and \(W_n\).  Thus \(V=W\) and the map is surjective.
This also proves closure of \(W\) under the stated Ihara bracket
without having to assume closure during the upper-bound argument.

Each homogeneous piece is finite-dimensional.  The isomorphism
therefore extends componentwise to the products of the weight
pieces.  The bracket in a fixed weight involves only finitely
many pairs of positive weights, so it extends to these products
and the two mutually inverse maps are continuous for the weight
filtration.  This proves the completed assertion.
\end{proof}
